\subsection*{T220\quad Debye声子谱的低温热容、零点能与高温极限}
三维各向同性晶体有$N$个原胞，声学声子在Debye近似中用截止频率$\omega_D$表征，三支声学模的总数为$3N$，定义$\Theta_D=\hbar\omega_D/k_B$。
(1) 从三维态密度的$\omega^2$标度确定Debye态密度和截止归一化。
(2) 计算声子热激发内能与零点能。
(3) 推导$T\ll\Theta_D$时的$T^3$热容定律及其系数。
(4) 检验$T\gg\Theta_D$时的Dulong–Petit极限并说明模型局限。
\par\textbf{逐问解析}
(1) 三维线性色散$\omega=c_sk$的波矢球壳态数正比$k^2dk$，因而声子态密度$g(\omega)=C\omega^2$。用$\int_0^{\omega_D}g(\omega)d\omega=3N$确定$C=9N/\omega_D^3$，所以$g(\omega)=9N\omega^2/\omega_D^3$。该归一化包含三支振动分支，而不是把模数错误地算为$N$。
(2) 每个玻色声子模热占有$\bar n_\omega=[e^{\beta\hbar\omega}-1]^{-1}$，故$U_{\rm th}=\int_0^{\omega_D}\hbar\omega\,\bar n_\omega g(\omega)d\omega=9Nk_BT(T/\Theta_D)^3\int_0^{\Theta_D/T}x^3/(e^x-1)\,dx$。零点能另为$U_{\rm zp}=\frac12\int_0^{\omega_D}\hbar\omega g(\omega)d\omega=\frac98Nk_B\Theta_D$，它与温度无关，不贡献定容热容。
(3) 在$T/\Theta_D\to0$时上限趋无穷，积分$\int_0^\infty x^3/(e^x-1)dx=\Gamma(4)\zeta(4)=\pi^4/15$。因此$U_{\rm th}\sim(3\pi^4/5)Nk_B T^4/\Theta_D^3$，定容微分得$C_V=(\partial U/\partial T)_V\sim(12\pi^4/5)Nk_B(T/\Theta_D)^3$。该指数来自三维低能态密度而非经验插值。
(4) 在$T\gg\Theta_D$，$e^x-1=x+O(x^2)$，热内能积分$\int_0^{\Theta_D/T}x^2dx=(\Theta_D/T)^3/3+O((\Theta_D/T)^4)$，于是$U_{\rm th}\sim3Nk_BT$、$C_V\to3Nk_B$。实际晶体的光学声子、各向异性和非谐相互作用会在Debye模型以外改变中高温行为及热膨胀。
\par\medskip
\subsection*{T221\quad 有简并度二能级系统的Schottky热容峰与最优温标}
$N$个彼此独立的粒子，各有基态能量0、简并度$g_0>0$与激发态能量$\Delta>0$、简并度$g_1>0$，温度$T$，无粒子数交换。
(1) 求单粒子配分函数和激发态占有几率。
(2) 推导热内能与定容热容的封闭表达。
(3) 对$g_0=g_1$求热容峰位置的超越方程及唯一正根。
(4) 分析低温和高温渐近，解释简并度对热容峰的影响。
\par\textbf{逐问解析}
(1) 以$x=\Delta/(k_BT)$计，单粒子配分函数$z=g_0+g_1e^{-x}$，全体系$Z_N=z^N$。定义$r=(g_1/g_0)e^{-x}$，激发态概率$p_1=g_1e^{-x}/z=r/(1+r)$，而基态概率$1-p_1$。简并度乘入Boltzmann权重，不能误把每一级当成一个微观态。
(2) 热内能$U=-\partial_\beta\ln Z_N=N\Delta p_1$，又$\partial_Tp_1=p_1(1-p_1)\Delta/(k_BT^2)$，因此$C_V=Nk_Bx^2r/(1+r)^2$。等价写成$C_V=N\operatorname{Var}(E)/(k_BT^2)$，因为$\operatorname{Var}(E)=\Delta^2p_1(1-p_1)$；这给出由正则涨落得到的第二条检查路线。
(3) 若$g_0=g_1$，有$r=e^{-x}$，$C_V/(Nk_B)=x^2/[4\cosh^2(x/2)]$。对$x>0$求导，驻点条件$2/x-\tanh(x/2)=0$，即$x\tanh(x/2)=2$。左边从0严格递增至无穷，故有唯一正根$x_*\simeq2.399357$，峰温$T_*=\Delta/(k_Bx_*)$；该热容峰是有限能级间热激发，不属于相变潜热。
(4) 低温$x\gg1$时$r\ll1$，$C_V\sim Nk_B(g_1/g_0)x^2e^{-x}$，指数冻结；高温$x\ll1$时$r\to g_1/g_0$，$C_V\sim Nk_Bx^2g_0g_1/(g_0+g_1)^2\to0$。简并度改变峰的高度和位置，但不会在有限能级的独立粒子模型中产生热力学奇异性。
\par\medskip
